% Esqueleto do artigo no formato EngUCP (Journal of Engineering of Catholic University of Petropolis).
% Compilar com PDFLaTeX (ex.: Overleaf): pdflatex, bibtex, pdflatex, pdflatex.
% Limites do modelo: 6 a 20 paginas (com figuras, tabelas e referencias), PDF <= 2 MB, sem numeracao de paginas.
% Orcamento de paginas indicado em cada secao. Numeros: copiar de redacao/01_resultados_unificado.md (cada um vem da planilha final).
\documentclass[oneside,a4paper,portuguese,links]{EngUCP}
\usepackage{graphicx}
\usepackage{amsmath,amsfonts}
\pagestyle{empty} % o modelo exige nao numerar as paginas

\title{O que as avaliações revelam sobre a nota: decomposição por aspecto do sentimento em avaliações de hóspedes do Airbnb no Rio de Janeiro}

\author[1]{Nome do Autor}      % <- ajustar nome, afiliacao e orientador
\affil[1]{Programa de Mestrado Profissional, Universidade Católica de Petrópolis, Petrópolis, RJ, Brasil}

\begin{document}
\vspace{3cm}
\maketitle

\let\oldthefootnote\thefootnote
\renewcommand{\thefootnote}{\fnsymbol{footnote}}
\footnotetext[1]{E-mail: \url{autor@ucp.br}.}
\let\thefootnote\oldthefootnote

\begin{keywords}
análise de sentimento por aspecto, avaliações online, Airbnb, modelagem de tópicos, diagnóstico de anúncios.
\end{keywords}

\begin{abstract}
% Rascunho (cabe na primeira pagina; sem figuras, tabelas, equacoes nem referencias).
As notas das plataformas de hospedagem são comprimidas no topo da escala e dizem pouco sobre o motivo da satisfação ou da insatisfação do hóspede.
Este trabalho usa o texto de 995.072 avaliações de anúncios do Airbnb no Rio de Janeiro (Inside Airbnb) para decompor a nota por aspecto e por queixa concreta.
O sentimento é medido por dois modelos multilíngues, e a escolha do modelo principal segue uma regra fixada antes da execução do segundo; os aspectos vêm de tópicos obtidos com BERTopic e são checados com um segundo método (LDA) e com um segundo anotador.
A medição é validada com frases rotuladas, ponderando a amostra pela população, e com as subnotas oficiais, e só limpeza, localização e custo-benefício atendem ao critério de confiança adotado.
Os resultados mostram que comodidades e limpeza concentram a maior parte da variação explicada da nota global, e que as queixas de instalações, limpeza e resposta do anfitrião são as mais associadas a notas menores, com ordem semelhante nos dois modelos.
Um índice sintético construído a partir do texto não acrescenta poder explicativo, além da nota oficial, sobre índices de volume de avaliações.
Os resultados são associações, e a aplicação proposta é um diagnóstico por anúncio, comparado ao bairro e ao tipo de acomodação.
\end{abstract}

% ---------------------------------------------------------------- (1,25 pag.)
\section{INTRODUÇÃO}
% Problema (nota saturada), pergunta central, quatro perguntas especificas, contribuicao, escopo (associacao, nao causa),
% organizacao do texto. Fonte: redacao/00_enquadramento.md.

% ---------------------------------------------------------------- (1,5 pag.)
\section{REFERENCIAL TEÓRICO}
% Avaliacoes online e compressao das notas; ABSA e modelos multilingues; BERTopic e LDA;
% validade e confiabilidade (concordancia entre anotadores, ponderacao por populacao); estudos sobre Airbnb. [escrever]

% ---------------------------------------------------------------- (3 pag.)
\section{DADOS E MÉTODO}
\subsection{Base e recorte}            % secao 1 do texto unificado
\subsection{O que a base permite medir} % secao 2: ocupacao/receita estimadas = indices de volume de avaliacoes
\subsection{Medição do sentimento e dos aspectos} % secao 3: modelos A e B, regra pre-registrada, BERTopic, mapa
\subsection{Desenho analítico}          % percentil da nota, LMG, GBM, logistica, CV 5x5, correcao BH
% Tabela 1 (opcional): resumo do corpus.

% ---------------------------------------------------------------- (3,5 pag.)
\section{VALIDAÇÃO DA MEDIÇÃO}
\subsection{Frases rotuladas e regra de escolha do modelo} % 4.1 e 4.3 (Tabela 2: acuracia e kappa, amostra e populacao)
\subsection{Segundo anotador e checagem com LDA}           % 4.5 e segundo anotador
\subsection{Classificação dos aspectos}                    % 4.6 (Tabela 3: confiavel x cautela)
% Concordancia A x B (4.2) e validade de construto detalhada (4.4): resumir em um paragrafo; tabelas completas na planilha.

% ---------------------------------------------------------------- (3 pag.)
\section{O QUE O TEXTO REVELA SOBRE A NOTA}
\subsection{Importância dos aspectos}   % 5.1 (Tabela 4: LMG A e B; Figura 1: barras da decomposicao)
\subsection{Queixas e elogios por tópico} % 5.2 (Tabela 5: 8 a 10 queixas de maior peso)
% Teste de efeito de tamanho de comodidades entra aqui, quando feito.

% ---------------------------------------------------------------- (1,25 pag.)
\section{APLICAÇÃO: DIAGNÓSTICO POR ANÚNCIO} % secao 6; uma figura com o relatorio de exemplo anonimizado

% ---------------------------------------------------------------- (1,75 pag.)
\section{UM ÍNDICE ÚNICO ACRESCENTA ALGO?}
% 7.1 e 7.2 em um paragrafo; Tabela 6: ganho de R2 fora da amostra (so amostra completa e nota = 5,0);
% resultado nulo com duas leituras (7.6). Coeficientes dentro da amostra e sensibilidades: planilha.

% ---------------------------------------------------------------- (1,25 pag.)
\section{DISCUSSÃO E LIMITAÇÕES} % secao 8 integrada a discussao com a literatura [escrever]

% ---------------------------------------------------------------- (0,75 pag.)
\section{CONCLUSÕES} % secao 9 + trabalhos futuros

\subsection*{\textit{Agradecimentos}}
\noindent Opcional.

\bibliographystyle{EngUCP}
\bibliography{referencias}
\end{document}
